% ============================================================
%  MTH636 – Game Theory | Lecture 23
%  Strategies and Equivalence in Imperfect Information Games
% ============================================================
\documentclass[12pt, a4paper]{article}

% ============================================================
%  MTH636 -- Game Theory
%  Notes Preamble — Color-Coded, TikZ-Enhanced
% ============================================================

% --- Core packages ---
\usepackage[T1]{fontenc}
\usepackage[utf8]{inputenc}
\usepackage{lmodern}
\usepackage[margin=2.5cm, top=2.8cm, bottom=2.8cm]{geometry}
\usepackage{microtype}

% --- Math ---
\usepackage{amsmath, amssymb, amsthm, mathtools}
\usepackage{bm}        % bold math
\usepackage{cancel}    % strike-through in math

% --- Color & Boxes ---
\usepackage[dvipsnames, svgnames, x11names, table]{xcolor}
\usepackage[most]{tcolorbox}
\tcbuselibrary{breakable, skins, theorems}

% --- TikZ & Diagrams ---
\usepackage{tikz}
\usepackage{tikz-qtree}
\usetikzlibrary{arrows.meta, positioning, shapes.geometric, fit, backgrounds,
                decorations.pathreplacing, calc, matrix, trees}
\usepackage{pgfplots}
\pgfplotsset{compat=1.18}

% --- Tables ---
\usepackage{booktabs}
\usepackage{multirow}
\usepackage{array}
\usepackage{tabularx}
\usepackage{makecell}

% --- Typography & Layout ---
\usepackage{enumitem}
\usepackage{parskip}
\usepackage{fancyhdr}
\usepackage{titlesec}
\usepackage{titletoc}
\usepackage{hyperref}
\usepackage{graphicx}
\usepackage{float}
\usepackage{caption}

% --- Color Palette ---
\definecolor{defblue}{HTML}{1A5276}      % Definition header
\definecolor{defbluebg}{HTML}{D6EAF8}    % Definition background
\definecolor{thmgreen}{HTML}{1E8449}     % Theorem header
\definecolor{thmbg}{HTML}{D5F5E3}        % Theorem background
\definecolor{exorange}{HTML}{B7770D}     % Example header
\definecolor{exbg}{HTML}{FEF9E7}         % Example background
\definecolor{warnred}{HTML}{922B21}      % Warning/Key Insight header
\definecolor{warnbg}{HTML}{FADBD8}       % Warning background
\definecolor{noticegray}{HTML}{2C3E50}   % Note header
\definecolor{noticebg}{HTML}{EBF5FB}     % Note background
\definecolor{taketawayyellow}{HTML}{7D6608} % Key Takeaway header
\definecolor{takeawaybg}{HTML}{FDFEFE}   % Takeaway background
\definecolor{accentpurple}{HTML}{6C3483} % Accent / proof
\definecolor{lightgray}{HTML}{F2F3F4}    % shading

% ============================================================
%  TCOLORBOX ENVIRONMENTS
% ============================================================

% -- Definition (Blue) --
\newtcolorbox{defbox}[2][]{
  enhanced, breakable,
  colback=defbluebg, colframe=defblue, coltitle=white,
  fonttitle=\bfseries\sffamily,
  title={Definition: #2},
  attach boxed title to top left={yshift=-2mm, xshift=4mm},
  boxed title style={colback=defblue, rounded corners},
  arc=4pt, boxrule=1.2pt,
  left=6pt, right=6pt, top=8pt, bottom=6pt,
  #1
}

% -- Theorem / Proposition (Green) --
\newtcolorbox{thmbox}[2][]{
  enhanced, breakable,
  colback=thmbg, colframe=thmgreen, coltitle=white,
  fonttitle=\bfseries\sffamily,
  title={Theorem: #2},
  attach boxed title to top left={yshift=-2mm, xshift=4mm},
  boxed title style={colback=thmgreen, rounded corners},
  arc=4pt, boxrule=1.2pt,
  left=6pt, right=6pt, top=8pt, bottom=6pt,
  #1
}

% -- Proposition (Green variant) --
\newtcolorbox{propbox}[2][]{
  enhanced, breakable,
  colback=thmbg, colframe=thmgreen!70!black, coltitle=white,
  fonttitle=\bfseries\sffamily,
  title={Proposition: #2},
  attach boxed title to top left={yshift=-2mm, xshift=4mm},
  boxed title style={colback=thmgreen!70!black, rounded corners},
  arc=4pt, boxrule=1.2pt,
  left=6pt, right=6pt, top=8pt, bottom=6pt,
  #1
}

% -- Example (Orange/Yellow) --
\newtcolorbox{exbox}[2][]{
  enhanced, breakable,
  colback=exbg, colframe=exorange, coltitle=white,
  fonttitle=\bfseries\sffamily,
  title={Example: #2},
  attach boxed title to top left={yshift=-2mm, xshift=4mm},
  boxed title style={colback=exorange, rounded corners},
  arc=4pt, boxrule=1.2pt,
  left=6pt, right=6pt, top=8pt, bottom=6pt,
  #1
}

% -- Key Insight / Warning (Red) --
\newtcolorbox{warnbox}[2][]{
  enhanced, breakable,
  colback=warnbg, colframe=warnred, coltitle=white,
  fonttitle=\bfseries\sffamily,
  title={Key Insight: #2},
  attach boxed title to top left={yshift=-2mm, xshift=4mm},
  boxed title style={colback=warnred, rounded corners},
  arc=4pt, boxrule=1.2pt,
  left=6pt, right=6pt, top=8pt, bottom=6pt,
  #1
}

% -- Note (Gray-Blue) --
\newtcolorbox{notebox}[1][]{
  enhanced, breakable,
  colback=noticebg, colframe=noticegray!60, coltitle=white,
  fonttitle=\bfseries\sffamily,
  title={Note},
  attach boxed title to top left={yshift=-2mm, xshift=4mm},
  boxed title style={colback=noticegray!60, rounded corners},
  arc=4pt, boxrule=0.8pt,
  left=6pt, right=6pt, top=8pt, bottom=6pt,
  #1
}

% -- Key Takeaway --
\newtcolorbox{takebox}[1][]{
  enhanced, breakable,
  colback=takeawaybg, colframe=taketawayyellow, coltitle=white,
  fonttitle=\bfseries\sffamily,
  title={Key Takeaway},
  attach boxed title to top left={yshift=-2mm, xshift=4mm},
  boxed title style={colback=taketawayyellow, rounded corners},
  arc=4pt, boxrule=1.2pt,
  left=6pt, right=6pt, top=8pt, bottom=6pt,
  #1
}

% -- Intuition --
\newtcolorbox{intubox}[1][]{
  enhanced, breakable,
  colback=accentpurple!8, colframe=accentpurple!60, coltitle=white,
  fonttitle=\bfseries\sffamily,
  title={Intuition},
  attach boxed title to top left={yshift=-2mm, xshift=4mm},
  boxed title style={colback=accentpurple!80, rounded corners},
  arc=4pt, boxrule=0.8pt,
  left=6pt, right=6pt, top=8pt, bottom=6pt,
  #1
}

% ============================================================
%  PAYOFF MATRIX HELPER
%  Usage: inside tabular, use \payoff{a}{b} for (a,b)
% ============================================================
\newcommand{\payoff}[2]{$(#1,\;#2)$}
\newcommand{\nash}[2]{\underline{$#1$},\underline{$#2$}}  % underline NE payoffs
\newcommand{\Nset}{\mathbb{N}}
\newcommand{\Rset}{\mathbb{R}}
\newcommand{\Sset}{\mathcal{S}}
\newcommand{\Ui}{u_i}

% ============================================================
%  SECTION FORMATTING
% ============================================================
\titleformat{\section}
  {\large\bfseries\sffamily\color{defblue}}
  {\thesection}{1em}{}[\titlerule]

\titleformat{\subsection}
  {\normalsize\bfseries\sffamily\color{thmgreen!80!black}}
  {\thesubsection}{1em}{}

\titleformat{\subsubsection}
  {\normalsize\itshape\sffamily\color{exorange!80!black}}
  {\thesubsubsection}{1em}{}

% ============================================================
%  HEADER / FOOTER
% ============================================================
\pagestyle{fancy}
\fancyhf{}
\fancyhead[L]{\sffamily\small\color{defblue} MTH636 -- Game Theory}
\fancyhead[R]{\sffamily\small\color{defblue} \leftmark}
\fancyfoot[C]{\sffamily\small\thepage}
\renewcommand{\headrulewidth}{0.4pt}

% ============================================================
%  HYPERREF SETUP
% ============================================================
\hypersetup{
  colorlinks=true,
  linkcolor=defblue,
  urlcolor=thmgreen,
  citecolor=accentpurple,
  bookmarksnumbered=true,
  pdftitle={MTH636 Game Theory Notes},
  pdfauthor={MTH636 Student Notes}
}

% ============================================================
%  MISC HELPERS
% ============================================================
\newcommand{\BR}{\text{BR}}        % Best Response
\newcommand{\NE}{\text{NE}}        % Nash Equilibrium
\newcommand{\SPNE}{\text{SPNE}}    % Subgame Perfect NE
\newcommand{\Dom}{\text{dom}}      % Dominance
\newcommand{\mixs}[1]{\sigma_{#1}} % Mixed strategy
\newcommand{\pures}[1]{s_{#1}}     % Pure strategy


% ============================================================
\begin{document}

% Title
\begin{center}
  {\LARGE\bfseries\sffamily\color{defblue} MTH636 -- Game Theory}\\[6pt]
  {\Large\sffamily Lecture 23: Strategies in Imperfect Information Games}\\[4pt]
\end{center}
\vspace{0.5cm}
\hrule height 1.5pt
\vspace{0.3cm}

% ============================================================
\section{Recap: The Formal Game}

We continue with the extensive form game from Lecture 22. Consider the game tree rooted at $n_0$ (Player I), where Player I chooses $A$, $B$, or $C$:
\begin{itemize}
  \item After $A$: reach $n_1$; after $B$: reach $n_2$. Player II has information set $h = \{n_1, n_2\}$, playing $\ell$ or $g$ from each. Outcomes: $n_4$ with $(4,1)$, $n_5$ with $(7,2)$ from $n_1$; $n_6$ with $(6,3)$, $n_7$ with $(4,0)$ from $n_2$.
  \item After $C$: reach $n_3$. Nature (O) moves, $\tfrac{1}{4}$ to $n_7$ (Player II moves $L$/$R$) and $\tfrac{3}{4}$ to $n_8$ (Player I moves $a$/$b$). After Player II's move at $n_7$: $n_{10}$ with $(0,0)$ and $n_{11}$ with $(8,4)$. After Player I's move at $n_8$: $n_{12}$ with $(8,8)$ and $n_{13}$ with $(0,8)$.
\end{itemize}

The information sets are:
\[
  V_1 = \{n_0, n_8\}, \quad V_2 = \{n_1, n_2, n_8\}
\]
\[
  H_1 = \bigl\{\{n_0\},\, \{n_8\}\bigr\}, \quad
  H_2 = \bigl\{\{n_1, n_2\},\, n_8\bigr\}
\]

\subsection{Actions at an Information Set}

For the information set $h = \{n_1, n_2\}$:
\[
  A(h) = \bigl\{\{\{n_1, n_4\},\, \{n_2, n_6\}\},\quad \{\{n_1, n_5\},\, \{n_2, n_7\}\}\bigr\}
\]
That is, the two actions at $h$ are $\ell$ (= choosing the left branch from whichever node in $h$ the player is at) and $g$ (= choosing the right branch), where each action bundles together one edge from each node in $h$.

% ============================================================
\section{Pure Strategies}

\begin{defbox}{Pure Strategy}
A \textbf{pure strategy} for player $i$, denoted $s_i$, is a mapping
\[
  s_i : H_i \;\longrightarrow\; \bigcup_{h \in H_i} A(h)
  \quad \text{such that} \quad
  s_i(h) \in A(h) \quad \forall\, h \in H_i.
\]
\end{defbox}

That is, a pure strategy assigns an action to every information set of player $i$, even those that may not be reached under the play induced by $s_i$.

The \textbf{pure strategy set} $S_i$ for player $i$ is the set of all such mappings.

The \textbf{pure strategy profile space} is:
\[
  S = S_1 \times S_2 \times \cdots \times S_n.
\]

The \textbf{utilities at strategy profiles} are defined as the \textbf{expected utilities} (since there may be chance nodes).

% ============================================================
\section{Mixed Strategies}

\begin{defbox}{Mixed Strategy}
\textbf{Mixed strategies} are probability distributions over pure strategies. For player $i$, the set of all mixed strategies is denoted $\Sigma_i$.
\end{defbox}

The \textbf{strategy profile} for mixed strategies is $\sigma = (\sigma_1, \ldots, \sigma_n)$.

\begin{notebox}
Utilities are defined as convex combinations of utilities at pure strategy profiles.
\end{notebox}

\subsection{Example}

Consider a two-player game tree: Player I ($n_0$) chooses $T_1$ or $B_1$ (reaching $n_1$ or $n_2$). Player II ($n_2$) has actions $L$ or $R$ (reaching $n_2$'s subtree). Player I also has a second information set containing $\{n_3, n_4\}$ where she plays $T_2$ or $B_2$. Player II has another information set $\{n_3, n_4, n_7\}$ where she plays $\beta_2$.

The pure strategy set of Player I is:
\[
  S_1 = \{T_1 T_2,\; T_1 B_2,\; B_1 T_2,\; B_1 B_2\}
\]

The mixed strategy set of Player I is:
\[
  \Sigma_1 = \Delta_4 = \left\{ (n_1, \ldots, n_4) \;\middle|\; n_i \geq 0,\; \sum_{i=1}^{4} n_i = 1 \right\}
\]

The set of \textbf{behaviour strategies} $B_1$ for Player I (explained below) is:
\[
  B_1 = \left\{ (n, 1-n),\, (y, 1-y) \;\middle|\; 0 \leq n \leq 1,\; 0 \leq y \leq 1 \right\}
\]
This is the set of behaviour strategies -- one independent randomisation per information set.

% ============================================================
\section{Behaviour Strategies}

\begin{defbox}{Behaviour Strategy}
A \textbf{behaviour strategy} for player $i$ is a mapping
\[
  b_i : H_i \;\longrightarrow\; \bigcup_{h \in H_i} \Delta\bigl(A(h)\bigr)
  \quad \text{such that} \quad
  b_i(h) \in \Delta\bigl(A(h)\bigr) \quad \forall\, h \in H_i.
\]
Here $\Delta(S)$ denotes a probability distribution over a set $S$.
\end{defbox}

Thus a behaviour strategy assigns to each information set an \emph{independent} probability distribution over the available actions at that information set.

\begin{itemize}
  \item For a behaviour strategy $b_i$, the probability that $b_i$ selects action $a$ at information set $h$ is denoted $b_i(a, h)$.
  \item A \textbf{behaviour strategy profile} is $b = (b_1, \ldots, b_n)$.
  \item With a slight abuse of notation, we write $\sigma = (\sigma_i)_{i \in N}$ for a strategy profile where $\sigma_i$ can be either a mixed or behaviour strategy for player $i$.
\end{itemize}

% ============================================================
\section{Reaching Probabilities}

\begin{defbox}{Reaching Probability}
For a mixed/behaviour strategy profile $\sigma = (\sigma_i)_{i \in N}$, and a node $n \in V$, we denote by $\rho(n; \sigma)$ the \textbf{probability that node $n$ will be reached} when the underlying strategy profile is $\sigma$.
\end{defbox}

\subsection{Example Computation}

Consider the following game tree: Player I ($n_0$) chooses $A$ or $B$ (reaching $n_1$ or $n_2$). Player II has information set $\{n_1\}$ (singleton), and plays $L$ or $R$ (reaching $n_3$ or $n_4$). Player I has a second information set $\{n_3, n_4\}$ where she plays $\ell$ or $r$ (reaching $n_5, n_6, n_7, n_8$).

Suppose:
\[
  \sigma_1 = \left[\tfrac{2}{10}(A\ell),\; \tfrac{1}{10}(Ar),\; \tfrac{4}{10}(B\ell),\; \tfrac{2}{10}(Br)\right]
  \quad \text{(mixed strategy for Player 1)}
\]
\[
  \sigma_2 = \left[\tfrac{1}{2}(L),\; \tfrac{1}{2}(R)\right]
  \quad \text{(mixed strategy for Player 2)}
\]

\paragraph{Computing $\rho(n_5; \sigma)$:}
$n_5$ will be reached when Player 1 plays $b$:
\[
  \rho(n_5; \sigma) = \tfrac{2}{10} \times \tfrac{1}{2} = \tfrac{2}{20} = \tfrac{1}{10}
\]

\paragraph{Computing $\rho(n_2; \sigma)$:}
$n_2$ will be reached if Player 1 plays either $A\ell$ or $Ar$, \emph{and} Player 2 plays $L$. Also the prob is:
\[
  \left[\tfrac{2}{10} + \tfrac{1}{10}\right] \times \tfrac{1}{2} = \tfrac{4}{20}
\]

% ============================================================
\section{Equivalence of Mixed and Behaviour Strategies}

\begin{defbox}{Equivalence of Mixed and Behaviour Strategies}
A \textbf{mixed strategy} $\sigma_i$ and a \textbf{behaviour strategy} $b_i$ of player $i$ are \textbf{equivalent} if for every mixed/behaviour strategy vector $\sigma_{-i}$ of the other players and every vertex $n \in V$,
\[
  \rho\bigl(n;\, (\sigma_i, \sigma_{-i})\bigr) \;=\; \rho\bigl(n;\, (b_i, \sigma_{-i})\bigr).
\]
\end{defbox}

That is, two strategies (one mixed, one behaviour) are equivalent if they induce identical probability distributions over the nodes of the game tree, regardless of what the other players do.

\begin{takebox}
\textbf{Key Insight.} This equivalence definition says that $\sigma_i$ and $b_i$ produce the same reaching probabilities for \emph{every} node in the game tree, for \emph{every} possible strategy profile of the opponents. Kuhn's theorem then states that under \textbf{perfect recall}, every mixed strategy has an equivalent behaviour strategy and vice versa.
\end{takebox}

\end{document}
